\chapter{Concurrency}
\label{ch:concurrency}

So far our programs have done one thing at a time, in order. \textbf{Concurrency}
is doing several things at once splitting work across the multiple cores every
modern computer has, so a big job finishes sooner. Salam expresses concurrency
with \emph{threads}, and protects shared data with \emph{synchronization}. This
chapter introduces both, along with the care they demand.

\section{Threads: running code in parallel}

A \textbf{thread} is an independent strand of execution. Your program already has
one the one running \code{main}. With \code{spawn} you can start another, which
runs a function of your choosing \emph{alongside} the rest of the program:

\begin{salam}
func greet():
    println "hello from a thread"
end

func main:
    handle : i64 = spawn(greet)
    println "hello from main"
    join(handle)
    println "the thread has finished"
end
\end{salam}

\code{spawn(greet)} starts \code{greet} running on a new thread and immediately
returns a \textbf{handle} an \code{i64} that identifies the thread while
\code{main} carries on. Because the two run concurrently, the order of the first
two lines of output is not guaranteed. \code{join(handle)} then \emph{waits} for
the thread to finish before continuing, so the final line always prints last.

\begin{note}
\code{spawn} takes a function with no parameters. To give a thread something to
work on, have it read shared variables (carefully see below) or use a closure
that captures the data it needs. \code{join} is how you know a thread is done; a
program should \code{join} every thread it spawns before relying on that thread's
results.
\end{note}

\section{The danger of shared data}

Threads become useful and dangerous the moment they touch the \emph{same}
data. Suppose two threads each add to a shared counter a thousand times. You would
expect the total to climb by two thousand. But incrementing a counter is really
three steps under the hood read the value, add one, write it back and if two
threads interleave those steps, they can read the same old value and each write
back the same new one, losing an increment. This is a \textbf{race condition}: the
result depends on the unpredictable timing of the threads, and it is one of the
hardest kinds of bug to find, because it appears only sometimes.

\section{Mutexes: one at a time}

The fix is to ensure that only one thread touches the shared data at a time. A
\textbf{mutex} (from ``mutual exclusion'') is a lock that exactly one thread can
hold. A thread \emph{locks} it before touching the shared data and \emph{unlocks}
it afterwards; any other thread that tries to lock it must wait until it is free.
The \code{sync} module provides them:

\begin{salam}
import "sync"

mut counter : int = 0
mut lock : sync.Mutex = sync.NewMutex()

func work():
    repeat 1000:
        sync.Lock(lock)
        counter = counter + 1     // protected: only one thread here at a time
        sync.Unlock(lock)
    end
end

func main:
    t1 : i64 = spawn(work)
    t2 : i64 = spawn(work)
    join(t1)
    join(t2)
    println "counter:", counter
    sync.Destroy(lock)
end
\end{salam}

\begin{progout}
counter: 2000
\end{progout}

Now the answer is reliably \code{2000}. Between \code{sync.Lock(lock)} and
\code{sync.Unlock(lock)} the \emph{critical section} only one thread ever
runs, so the read-add-write happens without interference. Create a mutex with
\code{sync.NewMutex()}, and call \code{sync.Destroy} on it when all the threads are
done.

\begin{warn}
A mutex protects data only if \emph{every} access to that data goes through it.
One thread that touches the shared counter without locking re-introduces the race.
And always pair a \code{Lock} with an \code{Unlock} on every path a lock that is
never released leaves other threads waiting forever (a \emph{deadlock}). Keep
critical sections short and make the lock/unlock pairing obvious.
\end{warn}

\section{Waiting for many threads}

When you spawn several threads and need to wait for all of them, you can collect
their handles and \code{join} each. For larger fan-outs, the \code{sync} module
also offers a \textbf{wait group} a counter you raise as you spawn work and
that threads decrement as they finish, letting the main thread block until the
count reaches zero:

\begin{salam}
import "sync"

mut wg : sync.WaitGroup = sync.NewWaitGroup()

func task():
    println "working"
    sync.Done(wg)         // signal completion
end

func main:
    sync.Add(wg, 3 as i64)
    spawn(task)
    spawn(task)
    spawn(task)
    sync.Wait(wg)         // block until all three call Done
    println "all tasks complete"
    sync.DestroyWaitGroup(wg)
end
\end{salam}

\code{sync.Add(wg, n)} announces \code{n} tasks; each task calls \code{sync.Done}
when finished; \code{sync.Wait} blocks until every \code{Done} has arrived. A wait
group is convenient when the number of workers is known but you do not want to
track each handle by hand.

\section{When concurrency pays and when it costs}

Concurrency is not free. Threads have a setup cost, locks add overhead, and ---
most of all concurrent code is harder to reason about and harder to debug than
sequential code. It pays off when there is genuine parallel work to do: a big
computation you can split into independent chunks, or many slow operations (such
as network requests) that can wait at the same time instead of one after another.

\begin{tip}
Reach for concurrency only when a measured need justifies it, and keep shared,
mutable state to an absolute minimum the less data threads share, the fewer
locks you need and the fewer races can occur. The simplest concurrent designs give
each thread its own data to work on and combine the results at the end, touching
shared state only briefly and under a lock. Start sequential; add threads when the
workload truly calls for them.
\end{tip}

\begin{deep}[In Depth: how threads map to the system]
Salam's \code{spawn} creates a real operating-system thread (a \code{pthread} on
Unix-like systems, a Windows thread on Windows), so threads run on separate cores
and your program can genuinely use a multi-core machine. The \code{sync.Mutex} is
a thin wrapper over the platform's native mutex. Because these are true OS
threads, the costs and the rules are the same as in C which is why the
discipline around shared data matters, and why the standard library gives you the
classic tools (mutexes, wait groups) to manage it.
\end{deep}

\section{Chapter summary}

\begin{itemize}[leftmargin=1.4em]
  \item \code{spawn(f)} runs the parameterless function \code{f} on a new thread
        and returns an \code{i64} handle; \code{join(handle)} waits for it to
        finish.
  \item Sharing mutable data between threads risks a \emph{race condition}, where
        the result depends on timing.
  \item A \code{sync.Mutex} (\code{NewMutex}, \code{Lock}, \code{Unlock},
        \code{Destroy}) ensures only one thread is in a critical section at a time.
  \item Every access to shared data must go through the lock, and every
        \code{Lock} must be matched by an \code{Unlock}.
  \item A \code{sync.WaitGroup} (\code{Add}, \code{Done}, \code{Wait}) waits for a
        known number of threads to complete.
  \item Use concurrency only when real parallel work justifies it, and minimise
        shared mutable state.
\end{itemize}

\section{Exercises}

\exercise Spawn a thread that prints the numbers 1 to 5 while \code{main} prints
the letters A to E. Join the thread, and observe that the interleaving varies
between runs.

\exercise Modify the shared-counter example to use \emph{four} threads of 1000
increments each and confirm the total is 4000.

\exercise Deliberately remove the \code{Lock}/\code{Unlock} from the counter
example, run it several times, and describe what you observe about the final total.

\exercise Use a \code{WaitGroup} to run five tasks that each print their own
number, and wait for all of them before printing ``done''.

\exercise Explain in a comment why keeping the critical section (the code between
\code{Lock} and \code{Unlock}) short is good practice.
